\documentclass[11pt,a4paper]{article}

\usepackage[margin=1in]{geometry}
\usepackage{amsmath,amssymb}
\usepackage{booktabs}
\usepackage{longtable}
\usepackage{array}
\usepackage[protrusion=true,expansion=false]{microtype}
\usepackage{xcolor}
\usepackage{titlesec}
\usepackage[hidelinks,colorlinks=true,linkcolor=linknavy,urlcolor=linknavy,citecolor=linknavy]{hyperref}
\usepackage{enumitem}
\usepackage{fancyhdr}
\usepackage{caption}
\usepackage{tabularx}
\usepackage[breakable]{tcolorbox}
\usepackage[numbers,sort&compress]{natbib}

\definecolor{linknavy}{HTML}{184A7A}
\definecolor{rulegray}{HTML}{9A9A94}
\definecolor{boxbg}{HTML}{F4F2EC}

\tcbset{
  keybox/.style={
    colback=boxbg, colframe=rulegray, boxrule=0.4pt, arc=2pt,
    left=8pt, right=8pt, top=6pt, bottom=6pt, breakable
  }
}

\titleformat{\section}{\large\bfseries}{\thesection}{0.6em}{}
\titleformat{\subsection}{\normalsize\bfseries}{\thesubsection}{0.6em}{}
\titlespacing*{\section}{0pt}{1.6ex plus .4ex}{0.9ex}

\newcommand{\ceff}{c_{\mathrm{eff}}}
\newcommand{\aSR}{\alpha_{SR}}
\newcommand{\aRS}{\alpha_{RS}}
\newcommand{\chat}{\hat{c}}

\setlength{\parskip}{0.55em}
\setlength{\parindent}{0pt}
\renewcommand{\arraystretch}{1.18}

\pagestyle{fancy}
\fancyhf{}
\fancyhead[L]{\small\itshape An identifiability audit of evolutionary oncology}
\fancyhead[R]{\small\thepage}
\renewcommand{\headrulewidth}{0.2pt}

\title{\vspace{-1.5em}\textbf{An Identifiability Audit of Evolutionary Oncology}\\[0.35em]
\large What standard-of-care measurement can and cannot recover ---\\and what it would take}

\author{Eric Schultz\\[0.2em]
\small Independent Researcher\\
\small ORCID 0009-0006-6283-1696}

\date{\small Working paper}

\begin{document}
\maketitle
\thispagestyle{fancy}

\begin{tcolorbox}[keybox]
\small
\textbf{Note on structure.} \textbf{Part I} (\S\S1--\ref{sec:limits}) states what we
believe. \textbf{Part II} (\S\ref{sec:errata}) states how we got there --- including
every hypothesis of our own that we had to abandon, and the one a referee had to find
for us.
\end{tcolorbox}

\section*{Abstract}

Adaptive and evolution-guided cancer therapies rest on quantitative claims about
tumour population dynamics. We ask a prior question: \textbf{given the data
oncology actually collects, can those claims be tested at all?}

They cannot. Fitting a two-population competition model to simulated patients
under uninterrupted maximum-tolerated-dose (MTD) therapy with q8w restaging,
\textbf{roughly half (50--57\%) of honest maximum-likelihood fits conclude that
resistance is free or advantageous when its true fitness cost is 30\%} --- worse
than a coin flip, with a median estimate of $\chat = -2.17$. The competition
coefficient is worse still: its estimate wanders the entire admissible range.
The Fisher information matrix has condition number ${\sim}10^{10}$ and the
likelihood is flat to the boundary. These parameters are not imprecise; they are
\emph{unidentified}.

\textbf{Denser sampling does not fix this.} A sevenfold increase in blood draws
leaves the sign-error rate exactly unchanged. Identifiability requires a
\emph{conjunction} of three things --- a clone-resolved observable, a window with
the drug \textbf{off}, and sampling concentrated around that window --- and
removing any one collapses it. Because treatment interruptions already occur in
routine care for clinical reasons, the triad can be assembled \textbf{without
prescribing any change in treatment}. Under this design the fitness cost resolves
at $n \approx 7$ patients and the directional suppression coefficient at
$n \approx 10$, at ${\sim}18$ blood draws each.

Crucially, \textbf{not every interruption is usable}. A hold for Grade 3/4 toxicity is
not a clean perturbation but a systemic crisis, and if the host insult spares the
resistant clone --- biologically plausible --- it biases the fitness cost
\emph{downward}, reproducing precisely the error the design exists to prevent
(sign-error returns to 58\%). The identifying set must be restricted to interruptions
that perturb the drug and nothing else: insurance lapses, supply interruptions,
scheduled surgery, non-adherence. The design can test this assumption on itself.

We then ask what the data \emph{should} be asked for, and find the field is asking
for the wrong thing. \textbf{The benefit of adaptive therapy is driven by
competition, not by the fitness cost of resistance.} A fitness cost is neither
necessary (36 days of benefit when resistance is \emph{free or advantageous}) nor
sufficient (4 days when the cost is large but competition is weak). Competition
carries roughly twice the leverage. This replicates in a mechanistic
consumer--resource model where competition \emph{emerges} from resource sharing
rather than being a free coefficient, and where a genuine zero-competition control
exists (leverage ratio 1.65 vs 1.95).

Finally, \textbf{scheduling requires no model at all}: closed-loop feedback
captures 86\% of an oracle's gain, where an open-loop schedule from a misspecified
model captures 6\%.

\vspace{1em}
\hrule
\vspace{0.6em}
\begin{center}\large\bfseries PART I --- WHAT WE BELIEVE\end{center}
\vspace{0.2em}
\hrule

\section{Introduction: a field with structural, not conceptual, deficits}

Cancer research does not suffer from a shortage of ideas or of rigour. It suffers,
like any large scientific enterprise, from deficits that are \emph{structural} ---
ones that no individual investigator is incentivised to repair, however clearly they
see them. We list them not as an indictment but to locate this paper among them, and
because the last item is the one we can actually do something about.

\begin{enumerate}[label=(\roman*),leftmargin=2.2em,itemsep=0.35em]
\item \textbf{Cancer is an evolving population, treated as a static target.}
Resistant subclones are near-inevitable in large heterogeneous metastatic
disease~\citep{gerlinger2012intratumor}, and MTD therapy plausibly
\emph{accelerates} their expansion by removing the sensitive competitors
suppressing them --- competitive release. Adaptive therapy has been formalised for
over a decade~\citep{gatenby2009adaptive}, validated preclinically
\citep{enriquez2016exploiting} and tested in metastatic castrate-resistant prostate
cancer~\citep{zhang2017integrating}; MTD remains standard of care.

\item \textbf{Metastasis and dormancy --- what actually kills patients --- are
studied less than primary tumour biology,} which is more tractable and more
publishable.

\item \textbf{Much of the foundational preclinical literature does not
replicate}~\citep{errington2021investigating,errington2021challenges,begley2012raise,prinz2011believe},
and no funded institution exists whose job is to confirm a preclinical result
before it consumes a phase III.

\item \textbf{Prevention and causation are underweighted relative to treatment,}
despite prevention accounting for several of the largest historical wins.

\item \textbf{Overdiagnosis and surrogate endpoints corrupt the measurement of
success.} Progression-free survival is not survival.

\item \textbf{Nobody runs the trials nobody profits from} --- off-patent agents,
generic combinations, de-escalation, schedule variation.

\item \textbf{The measurement protocol cannot identify the quantities the models
depend on.}
\end{enumerate}

This paper is about (vii). It is \textbf{not} a claim that adaptive therapy works.
It is a narrower and more defensible claim: \textbf{under current measurement
practice, adaptive therapy's central premises are not testable} --- and this is
fixable, cheaply, without asking any clinician to treat any patient differently.

\section{The two candidate parameters}

Adaptive therapy's mechanism is \emph{competitive suppression}: keep drug-sensitive
cells alive so they hold resistant cells in check. Two parameters could underwrite
this, and the field has fixated on one of them.

\textbf{The fitness cost of resistance, $c$.} Resistant cells are assumed to grow
more slowly without drug: $r_R = r_S(1-c)$, a life-history trade-off
\citep{aktipis2013life} with direct empirical support in EGFR-mutant lung cancer
\citep{chmielecki2011optimization}. The standard framing is that adaptive therapy
works \emph{if and only if} $c > 0$~\citep{gatenby2009adaptive}. This is the
parameter the literature argues about.

\textbf{The competition coefficient, $\alpha$.} Specifically $\aSR$: the degree to
which sensitive biomass denies carrying capacity to resistant cells. There is direct
experimental evidence that spatial competition alone constrains resistance
\citep{bacevic2017spatial}. This is the parameter nobody measures.

We will show that neither is identifiable under standard of care; that both become
identifiable under a cheap observational redesign; and that \textbf{$\alpha$, not
$c$, is the one that decides the question.}

\section{Model and audit design}

\textbf{Dynamics.} Two competing populations under therapy:
\begin{align}
\frac{dS}{dt} &= r_S\, S \left(1 - (S + \aRS R)\right) - d\, u(t)\, S\\
\frac{dR}{dt} &= r_R\, R \left(1 - (R + \aSR S)\right), \qquad r_R = r_S(1-c)
\end{align}
with $u(t) \in \{0,1\}$ the drug indicator. Truth: $r_S = 0.03$/d, $c = 0.30$,
$\alpha = 1.0$, $d = 0.05$/d, initial resistant fraction $f_0 = 0.01$.

\textbf{Observation.} ctDNA is cleared from plasma with a half-life of minutes to
a few hours~\citep{lo1999rapid,diehl2008circulating,esmo2022ctdna}, so plasma
concentration is a near-instantaneous readout of current shedding rather than a
smoothed average. Shedding has a burden-proportional and a death-proportional
component, the latter reproducing the transient ctDNA \emph{rise} seen days after
cytotoxic initiation~\citep{esmo2022ctdna}:
\begin{equation}
y(t) = (S + \rho R) + \kappa\, d\, u(t)\, S
\end{equation}

Two observables are compared:
\begin{itemize}[leftmargin=1.4em,itemsep=0.2em]
\item \textbf{total} --- plasma ctDNA burden alone (a tumour-agnostic assay), 30\%
CV lognormal noise.
\item \textbf{clone-resolved} --- burden \textbf{and} the resistant allele fraction
(a tumour-informed panel), on the logit scale.
\end{itemize}

\textbf{Protocols.} \emph{Continuous} MTD (standard of care); and an
\emph{incidental hold} --- an unplanned interruption of $L$ days, of the kind that
occurs routinely for toxicity, dose reduction, surgery, insurance gaps, or
non-adherence.

\textbf{Inference.} We distinguish \emph{structural} from \emph{practical}
non-identifiability throughout~\citep{raue2009structural,miao2011identifiability},
following the approach of \citet{eisenberg2017confidence} for cancer chemotherapy
models. We use Fisher information and Cram\'er--Rao bounds, validated throughout
against \textbf{multi-start} maximum likelihood on synthetic patients,
with $c$ free on $[-3, 0.99]$ so a fit is \emph{permitted} to conclude resistance
is beneficial. Restarts are drawn from the whole parameter box, not perturbed from
the truth.

\textbf{The test.} Not ``is $c$ estimated precisely'' but the weakest possible bar:
\textbf{can we determine its sign?}

\section{Under standard of care, nothing is identifiable}

\subsection{The fitness cost}

\begin{center}
\begin{tabular}{lrrrc}
\toprule
\textbf{Sampling} & \textbf{Draws} & \textbf{cond(FIM)} & \textbf{SE($c$)} & \textbf{Sign resolved?}\\
\midrule
\textbf{standard (q8w)} & 6 & $6.6\times10^{10}$ & 5.21 & no\\
q4w & 11 & $5.8\times10^{10}$ & 3.72 & no\\
q2w & 21 & $6.3\times10^{10}$ & 3.03 & no\\
weekly & 41 & $6.8\times10^{10}$ & 2.44 & no\\
\textbf{daily} & 281 & $7.6\times10^{10}$ & 1.07 & \textbf{no}\\
\bottomrule
\end{tabular}
\end{center}

\emph{Total ctDNA, continuous MTD.} \textbf{Daily sampling for forty weeks --- 281
draws --- fails.} The condition number is ${\sim}10^{10}$ throughout and barely
moves with cadence: adding measurements does not add information about $c$.

The estimator correlation matrix shows why: $\chat$ correlates with $\hat{r}_S$ at
0.996, and $\hat{\alpha}$ with $\hat{f}_0$ at 1.000. \textbf{Total burden cannot
distinguish a slow-growing resistant clone from a fast-growing clone that started
rarer.} The clonal split is not in the observable.

At these condition numbers a Cram\'er--Rao bound is a linearisation of a
near-singular surface, so we fit synthetic patients directly. This degradation with
parameter count is consistent with what is already known for simpler growth models:
two-parameter models are practically identifiable while three-parameter models
frequently are not, even on dense preclinical
data~\citep{simpson2022parameter,porthiyas2024practical}, and lack of practical
identifiability is a known route to confidently wrong
predictions~\citep{gallo2022lack}.

\subsection{Multi-start MLE: worse than the bound suggests}

\begin{center}
\small
\begin{tabular}{lccc}
\toprule
\textbf{Design} & \textbf{median $\chat$} & \textbf{95\% interval} & \shortstack{\textbf{fits concluding resistance}\\\textbf{is FREE or BENEFICIAL}}\\
\midrule
\textbf{status quo} --- total, q8w, MTD & $\mathbf{-2.17}$ & $[-3.00, +0.88]$ & \textbf{50--57\%}\\
dense only --- total, \textbf{weekly}, MTD & $-0.34$ & $[-3.00, +0.86]$ & \textbf{57\%}\\
clones only --- cloneAF, q8w, MTD & $+0.25$ & $[-3.00, +0.74]$ & \textbf{37\%}\\
\textbf{proposed} --- cloneAF, burst, 28d hold & $\mathbf{+0.33}$ & $[-0.06, +0.51]$ & \textbf{7--10\%}\\
\bottomrule
\end{tabular}
\end{center}

The lower bound hits the \emph{box constraint} in the first three rows: the
likelihood is flat to the boundary. The parameter is unidentified, and the ``SE''
of 5.21 is a meaningless number attached to a flat surface.

\begin{tcolorbox}[keybox]
In a world where resistance carries a substantial 30\% fitness cost, a competent
analyst fitting a competent model to competently collected data concludes that
\textbf{resistance is free or advantageous roughly half the time}. The median
honest estimate is $\chat = -2.17$: confidently, wildly, in the wrong direction.
\end{tcolorbox}

\subsection{The competition coefficient is worse}

\begin{center}
\small
\begin{tabular}{lccc}
\toprule
\textbf{Design} & \textbf{median $\hat{\alpha}$} & \textbf{95\% interval} & \textbf{$\alpha=0$ excluded?}\\
\midrule
status quo & 3.000 \emph{(pinned at bound)} & $[0.000, 3.000]$ --- \textbf{the entire box} & no\\
\textbf{proposed} & \textbf{0.867} & $\mathbf{[0.471,\ 1.573]}$ & \textbf{yes}\\
\bottomrule
\end{tabular}
\end{center}

Under standard of care the estimate for $\alpha$ wanders the whole admissible
range. There is essentially no information in the data about it.

\section{The identifiability triad}

\textbf{The three ingredients are conjunctive, not additive.} Removing any one
collapses identifiability.

\begin{itemize}[leftmargin=1.4em,itemsep=0.3em]
\item \textbf{Sampling density alone:} 6 draws $\to$ 41 draws (a $7\times$ increase
in phlebotomy). Sign-error: 57\% $\to$ \textbf{57\%}. \emph{Exactly unchanged.}
\item \textbf{Clone resolution alone:} total ctDNA $\to$ total $+$ resistant allele
fraction, at q8w. Sign-error: 57\% $\to$ \textbf{37\%}. Better; still catastrophic.
\item \textbf{Both, plus a 28-day drug-off window:} \textbf{7--10\%}.
\end{itemize}

\subsection{Why the drug-off window is indispensable}

Under continuous MTD the sensitive population is driven toward extinction, the
competition term $\alpha S$ vanishes, and the resistant clone expands essentially
unopposed at $r_R$ --- which is confounded with $r_S$ and $f_0$.

\begin{tcolorbox}[keybox]
A drug-off window is the only thing that lets the sensitive population re-express
its growth advantage --- which is what competitive suppression \emph{is}. This is a
system-identification truism: \textbf{you cannot identify a system you never
excite.} Evolutionary oncology has been trying to.
\end{tcolorbox}

\subsection{What the interruption must look like}

\begin{center}
\begin{tabular}{lcc}
\toprule
\textbf{Hold length} & \textbf{SE($c$)} & \textbf{Sign resolved (single patient)?}\\
\midrule
7 d & 0.394 & no\\
21 d & 0.191 & no\\
\textbf{28 d} & \textbf{0.151} & \textbf{yes}\\
42 d & 0.107 & yes\\
\bottomrule
\end{tabular}
\end{center}

Hold \emph{timing} matters far less than hold \emph{length} (SE 0.177--0.247 for
holds starting anywhere between days 28 and 196). This is fortunate: one cannot
choose when a toxicity hold occurs, but one can choose to sample around it.

\subsection{Design efficiency, in patients}

Patients required to determine the sign of a subtle cost ($c = 0.15$), ideal-pooling
bound:

\begin{center}
\begin{tabular}{lllrr}
\toprule
\textbf{Observable} & \textbf{Sampling} & \textbf{Protocol} & \textbf{Draws/pt} & \textbf{Patients}\\
\midrule
total & q8w & uninterrupted & 6 & \textbf{89{,}424}$^\dagger$\\
total & burst & uninterrupted & 18 & 23{,}087$^\dagger$\\
total & burst & 28d hold & 18 & 150\\
cloneAF & q8w & uninterrupted & 6 & 102\\
cloneAF & q8w & 28d hold & 6 & 23\\
\textbf{cloneAF} & \textbf{burst} & \textbf{28d hold} & \textbf{18} & \textbf{7}\\
\bottomrule
\end{tabular}
\end{center}

{\small $^\dagger$ Read as ``unidentified'', not as a literal sample size --- these
are extrapolations from a singular surface.}

The \textbf{directional} suppression coefficient $\aSR$ (with $\aSR$ and $\aRS$ kept
separate --- the hard case) resolves at $\mathbf{n \approx 10}$:

\begin{center}
\begin{tabular}{cccc}
\toprule
\textbf{patients} & \textbf{median $\hat{\alpha}_{SR}$} & \textbf{95\% interval} & \textbf{resolved?}\\
\midrule
1 & 1.201 & $[-1.000, 3.000]$ & no\\
5 & 0.774 & $[-0.733, 2.184]$ & no\\
\textbf{10} & \textbf{1.048} & $\mathbf{[0.472,\ 1.941]}$ & \textbf{yes}\\
40 & 0.986 & $[0.511, 1.322]$ & yes\\
\bottomrule
\end{tabular}
\end{center}

\textbf{No single patient's data will do it; a cohort of ten will.} This is a
\emph{cohort} design, not a personalised one --- it can establish whether
competitive suppression is real in a tumour type, not whether it is operating in
the patient in front of you.

\section{Robustness, and the right estimand}

\subsection{What the design absorbs}

\begin{itemize}[leftmargin=1.4em,itemsep=0.3em]
\item \textbf{Differential clone shedding.} If clones shed ctDNA at different rates
($\rho = \beta_R/\beta_S \neq 1$), the plasma allele fraction is a biased view of
clonal composition. Making $\rho$ a free parameter costs \textbf{nothing}: SE($c$)
unchanged at 0.151. Even \emph{wrongly} assuming $\rho = 1$ against a threefold
truth biases $\chat$ by only $-0.037$.
\item \textbf{Shedding-model ($\kappa$) misspecification.} A $2.5\times$ error
biases $\chat$ by $-0.02$ to $-0.05$. Conservative in direction.
\item \textbf{Between-patient heterogeneity.} Inflation is mild while
$\omega < \mathrm{SE}_1$ ($n$: 10 $\to$ 14 at $c = 0.15$ with $\omega = 0.15$).
\item \textbf{Competition regime.} SE($c$) $=0.151$ for the proposed design at every
$\alpha$ from 0.5 to 1.0. Identifiability is regime-robust.
\end{itemize}

\subsection{What it does not absorb: a fixed compartment count}

Real resistance is polyclonal and spatially structured
\citep{gerlinger2012intratumor,diaz2012molecular,misale2012emergence}. Fitting a
two-clone model to three-clone truth destroys the estimate --- \textbf{and the sign
of the bias is set by modelling convention, not by data:}

\begin{center}
\begin{tabular}{lcc}
\toprule
\textbf{2-clone model variant} & \textbf{median $\chat$} & \textbf{bias}\\
\midrule
rigid kill (resistant assumed fully refractory) & 0.369 & $\mathbf{+0.069}$\\
$+$ differential shedding & 0.445 & $\mathbf{+0.145}$\\
\textbf{flexible kill} (kill rate free) & $\mathbf{-0.12}$ to $\mathbf{-0.04}$ & $\mathbf{\approx -0.4}$, \textbf{sign flips}\\
\bottomrule
\end{tabular}
\end{center}

Same data, same clone count, opposite conclusion. \textbf{Pooling divides variance
by $n$ and does nothing to bias.} At large $n$ one converges, with tight confidence
intervals, on a number whose sign was decided by a modelling choice.

\subsection{The ontology is undecidable --- and does not need to be decided}

Within a \emph{discrete} candidate set, clone number is recoverable \textbf{only
with a per-variant panel}:

\begin{center}
\begin{tabular}{llc}
\toprule
\textbf{Truth} & \textbf{Assay} & \textbf{$k$ recovered}\\
\midrule
$k=3$ & \textbf{lumped} allele fraction & \textbf{0\%} --- selects $k{=}2$ in 10/10\\
$k=3$ & \textbf{per-variant panel} & \textbf{100\%}\\
\bottomrule
\end{tabular}
\end{center}

The lumped assay reports a plausible-looking $\chat = 0.185$ with no diagnostic
complaint. Silent, confident, unfalsifiable from within the analysis.

But admit a \textbf{continuum} model --- resistance as a distribution of
sensitivities rather than compartments --- and clone number evaporates:

\begin{center}
\begin{tabular}{lcccc}
\toprule
\textbf{Truth} & \textbf{$\hat{k}=2$} & \textbf{$\hat{k}=3$} & \textbf{continuum} & \textbf{correct?}\\
\midrule
CONTINUUM & 1/8 & 0/8 & \textbf{7/8} & yes\\
DISCRETE-3 & 2/8 & 0/8 & \textbf{6/8} & \textbf{no --- spurious}\\
\bottomrule
\end{tabular}
\end{center}

\textbf{BIC prefers the continuum in both worlds.} The compartment count is not
identifiable.

\textbf{It does not need to be.} The clone number was never the estimand. What
adaptive therapy requires is a \emph{functional} of the resistance distribution,
well-defined under either ontology:
\begin{equation}
\ceff = 1 - \frac{\mathbb{E}\!\left[\, r(x) \mid x > 0 \,\right]}{r_S}
\end{equation}

\begin{center}
\begin{tabular}{lccc}
\toprule
\textbf{Truth} & \textbf{true $\ceff$} & \textbf{reported $\hat{c}_{\mathrm{eff}}$} & \textbf{wrong sign}\\
\midrule
CONTINUUM & 0.121 & 0.075 & \textbf{0\%}\\
DISCRETE-3 & 0.200 & 0.119 & \textbf{0\%}\\
\bottomrule
\end{tabular}
\end{center}

\textbf{The sign is recovered in both worlds, under the wrong ontology, every
time} --- with a bias that \emph{understates} the cost, erring against the
hypothesis being tested.

Two requirements are non-negotiable: a \textbf{per-variant panel} (a lumped allele
fraction recovers nothing, and reports a plausible wrong answer silently), and a
\textbf{continuum model in the candidate set} (discrete-only candidate sets inflate
$\hat{c}_{\mathrm{eff}}$ by $2$--$2.6\times$, in the direction that flatters
adaptive therapy).

\begin{tcolorbox}[keybox]
Any published fitness-cost estimate derived from a small fixed number of
compartments fitted to a lumped allele fraction should be treated as
\textbf{uninterpretable} --- not as an estimate with wide error bars, but as a
number whose sign was determined by a modelling convention.
\end{tcolorbox}


\subsection{The host is not a neutral bystander: which interruptions are actually clean}

An earlier draft of this paper treated \emph{any} incidental treatment interruption
as a free excitation. That was wrong, and the error was material.

A hold for Grade 3/4 toxicity is not a clean drug-off window. The patient is in
systemic crisis --- hepatotoxicity, neutropenia, inflammation --- and the growth and
competition parameters of a tumour in a metabolically deranged host are not those of
a tumour in a well one. We model this explicitly. During the hold, let
$r_i \to r_i(1 + \delta_i)$ and $\alpha \to \alpha(1 + \delta_\alpha)$, generate data
from the \emph{sick} host, and fit the naive model that assumes the window is clean.

\begin{center}
\small
\begin{tabular}{lccc}
\toprule
\textbf{Scenario} & \textbf{median $\chat$} & \textbf{bias} & \textbf{sign wrong}\\
\midrule
\textbf{clean (non-toxic) hold} & $+0.301$ & $+0.001$ & \textbf{9\%}\\
symmetric insult (both clones $\times 0.60$) & $+0.238$ & $-0.062$ & 29\%\\
\textbf{differential} (sens $\times0.60$, res $\times0.80$) & $+0.146$ & $-0.154$ & \textbf{36\%}\\
\quad \emph{same data, host effect fitted as nuisance} & $+0.241$ & $-0.059$ & 22\%\\
\textbf{severe differential} (sens $\times0.40$, res $\times0.90$) & $\mathbf{-0.178}$ & $\mathbf{-0.478}$ & \textbf{58\%}\\
\bottomrule
\end{tabular}
\end{center}

Three findings, and the structure of them matters more than the numbers.

\textbf{(1) A \emph{symmetric} host insult is comparatively benign.} Because $c$ is
defined by a \emph{ratio}, $r_R = r_S(1-c)$, a factor common to both clones largely
cancels. Severe symmetric slowing still degrades the estimate --- not by biasing it
but by \emph{attenuating the excitation itself}, since the information in a drug-off
window scales with how fast the clones diverge during it --- but it does not invert
the sign.

\textbf{(2) A \emph{differential} insult that spares the resistant clone is fatal.}
This does not cancel; it is directly confounded with $c$. And it is the biologically
plausible case: efflux pumps, stress-response programs and quiescence plausibly
confer \emph{both} drug resistance \emph{and} robustness to host metabolic crisis.
Worse, it biases $\chat$ \textbf{downward}.

\begin{tcolorbox}[keybox]
\textbf{The confound is aligned with the failure mode.} A host insult that spares the
resistant clone makes resistance \emph{look} free --- which is precisely the false
conclusion this paper exists to prevent. It does not add noise; it manufactures our
error, and it does so with a straight face.
\end{tcolorbox}

\textbf{(3) Fitting the host effect as a nuisance parameter does not rescue it}
(36\% $\to$ 22\%). We tried the obvious statistical fix and it is insufficient.

\textbf{The remedy is therefore not statistical but an inclusion criterion.} Not every
incidental hold carries a host insult. Interruptions arising from
\emph{insurance lapse or prior-authorisation delay, drug supply interruption,
scheduled surgery, elective procedure, travel, non-adherence, or clinic closure}
perturb the drug and nothing else. \textbf{These, and only these, are clean
excitations.} Toxicity holds must be excluded from the identifying set --- or, at
minimum, analysed as a separate stratum and never pooled with the clean ones.

The cost is enrolment, not validity: if a fraction $f$ of interruptions are non-toxic,
enrol $n/f$ patients to obtain $n$ usable ones. At a plausible $f \approx 0.3$--$0.5$,
that is roughly 15--35 patients enrolled to yield the 7--10 the design requires. Still
a small study.

This tightens the design rather than weakening it, and it yields a new falsifiable
prediction (\S\ref{sec:predict}): \emph{within the same cohort, $\hat{c}$ estimated from toxic
holds should be systematically lower than $\hat{c}$ estimated from non-toxic holds.}
If the two strata agree, the confound is absent. If they diverge, the design has
detected its own principal threat --- from the inside, using data it was already
collecting.


\section{What actually drives adaptive therapy}

\subsection{Competition, not cost}

Sweeping the fitness cost $c$ (including \textbf{negative} values, where resistance
is advantageous) against the competition coefficient $\alpha$, restricted to the
\textbf{stable-coexistence regime} ($\alpha < 1$ --- above that the populations
cannot coexist at all, which is the thing adaptive therapy depends on).

\textbf{Absolute gain (days) of model-free adaptive therapy over MTD:}

\begin{center}
\begin{tabular}{lrrrr}
\toprule
$\alpha \downarrow$ \textbf{/} $c \rightarrow$ & $-0.20$ & $0.00$ & $+0.20$ & $+0.40$\\
\midrule
\textbf{0.10} & $-1$ & 0 & 1 & \textbf{4}\\
0.40 & 8 & 12 & 17 & 27\\
0.70 & 16 & 30 & 33 & 51\\
\textbf{0.85} & \textbf{30} & \textbf{36} & 39 & \textbf{66}\\
\bottomrule
\end{tabular}
\end{center}

Leverage of $\alpha$: \textbf{43 days.} Leverage of $c$: \textbf{22 days.}
\textbf{Ratio $\approx 2:1$.}

\begin{enumerate}[leftmargin=1.6em,itemsep=0.3em]
\item \textbf{A fitness cost is not necessary.} 36 days of benefit at $c \leq 0$,
where resistance is free or \emph{advantageous}.
\item \textbf{A fitness cost is not sufficient.} At weak competition, a \emph{large}
cost ($c = +0.40$) buys \textbf{4 days}.
\item \textbf{$\alpha$ carries roughly twice the leverage.}
\end{enumerate}

The field's ``if and only if'' is false in both directions.

\begin{tcolorbox}[keybox]
Sensitive cells suppress resistant ones primarily by \textbf{occupying space they
need}, not by \textbf{outgrowing} them. A fitness cost amplifies the effect; it does
not create it.
\end{tcolorbox}

This is not without empirical precedent. \citet{bacevic2017spatial} report that
\emph{spatial} competition alone constrains resistance to targeted therapy --- a
result our audit predicts, and which sits awkwardly with a literature that continues
to frame adaptive therapy as contingent on a fitness cost.

\subsection{This is not an artefact of Lotka--Volterra}

In LV, $\alpha$ and $c$ are independent free parameters --- and that independence is
what makes the comparison possible. So we rebuilt the test in \textbf{MacArthur's consumer--resource
model}~\citep{macarthur1970species}, the mechanistic model from which LV competition
coefficients are \emph{derived}, with $\alpha \propto$ niche overlap $\omega$:
\begin{align}
\mathrm{uptake}_S &= a_S\,(\mathrm{Res}_1 + \omega\,\mathrm{Res}_2)\\
\mathrm{uptake}_R &= a_R\,(\omega\,\mathrm{Res}_1 + \mathrm{Res}_2), \qquad a_R = a_S(1-c)
\end{align}

Here competition \textbf{emerges} from resource sharing rather than being assumed;
$\omega = 0$ is a \textbf{genuine zero-competition control} (the clones eat
different resources); and $c$ is a separate trait (uptake efficiency), not collapsed
into competitive ability.

\begin{center}
\begin{tabular}{llrrr}
\toprule
$\omega$ & \textbf{competition} & $c=-0.10$ & $0.00$ & $+0.30$\\
\midrule
\textbf{0.00} & \textbf{NONE} & $\mathbf{-0}$ & $\mathbf{-0}$ & $\mathbf{-0}$\\
0.50 & moderate & 96 & 107 & 240\\
1.00 & maximal & 124 & 147 & 298\\
\bottomrule
\end{tabular}
\end{center}

\begin{center}
\begin{tabular}{lcc}
\toprule
 & \textbf{Lotka--Volterra} & \textbf{MacArthur}\\
\midrule
\textbf{ratio competition : cost} & \textbf{1.95} & \textbf{1.65}\\
benefit with \textbf{no competition}, best $c$ & $\approx 0$ d & \textbf{exactly 0 d}\\
benefit with \textbf{no fitness cost} & 36 d & \textbf{148 d}\\
\bottomrule
\end{tabular}
\end{center}

\textbf{The result replicates.} Different model, different mechanism, a real
control --- the same conclusion.

\subsection{Scheduling needs no model}

Adaptive therapy's clinical rule is \textbf{closed-loop}: treat until burden falls
to 50\% of baseline, stop, resume at baseline --- the rule used in the metastatic
prostate cancer trial~\citep{zhang2017integrating}. It has \emph{no parameters}.

\begin{center}
\begin{tabular}{lrrrr}
\toprule
\textbf{Truth} & \textbf{MTD} & \textbf{Gatenby (no model)} & \textbf{Open-loop (misspec.)} & \textbf{Oracle}\\
\midrule
continuum & 637 d & \textbf{847 d} & \textbf{651 d} & 882 d\\
\bottomrule
\end{tabular}
\end{center}

The model-free feedback rule captures \textbf{86\%} of the oracle's available gain.
The open-loop schedule derived from a misspecified model captures \textbf{6\%} ---
while looking principled.

\begin{tcolorbox}[keybox]
\textbf{Feedback substitutes for identification --- in \emph{execution}, not in
\emph{selection}.} You do not need an accurate model to \emph{run} the therapy, given
an accurate sensor and a stabilising feedback law. You still need $\alpha$ to know
\emph{whom to run it on}.
\end{tcolorbox}

This distinction is easy to miss and it does not undercut the trial --- it is the
trial's whole justification. A feedback controller is only worth deploying in a tumour
where competitive suppression actually operates; at $\omega = 0$ or $\alpha \approx 0$
the loop captures nothing (\S7.1--7.2), and the patient has been managed at
sub-maximal dose for no benefit. \textbf{Identification is what tells you which
cancers, and which patients, are candidates for the loop in the first place.} The
measurement is for \emph{selection}; the feedback is for \emph{execution}. Neither
replaces the other.

\section{The proposed design}

The audit demands a drug-off window, which looks like an intervention. It is not ---
because \textbf{the perturbation already exists.}

Treatment interruptions occur constantly in routine oncology for reasons that have
nothing to do with evolutionary theory. From a system-identification standpoint these
are \textbf{free excitations of the system}, and they are currently wasted, because
nobody samples around them.

But --- and this is the correction forced by \S6.4 --- \textbf{not all of them are
clean}. The identifying set must be restricted to interruptions that perturb
\emph{the drug and nothing else}:

\begin{center}
\small
\begin{tabular}{ll}
\toprule
\textbf{Clean (use these)} & \textbf{Confounded (exclude, or stratify separately)}\\
\midrule
insurance lapse / prior-auth delay & Grade 3/4 toxicity hold\\
drug supply interruption & hepatotoxicity, neutropenia, sepsis\\
scheduled surgery or elective procedure & severe fatigue / performance-status collapse\\
travel, clinic closure, holiday & intercurrent serious illness\\
non-adherence, physician discretion & dose reduction for tolerability\\
\bottomrule
\end{tabular}
\end{center}

A hold caused by an insurance company tells you nothing about the patient's
physiology. A hold caused by Grade 4 neutropenia tells you a great deal, and all of it
is confounded with what you are trying to measure.

\begin{tcolorbox}[keybox]
\textbf{We do not ask you to stop the drug.} We ask that when you stop it --- for
whatever clinical reason you judge appropriate --- we may draw blood densely around
that window, and that the assay resolve clones rather than only total burden.
\end{tcolorbox}

Concretely:

\begin{enumerate}[label=(\alph*),leftmargin=1.8em,itemsep=0.3em]
\item A tumour-informed panel reporting \textbf{each resistance variant separately}
--- not tumour-agnostic total ctDNA, and not a single merged allele fraction.
Serial ctDNA tracking of specific resistance variants is already established
practice~\citep{dawson2013analysis,diaz2012molecular,misale2012emergence}; what is
missing is that the variants be reported separately rather than merged. This
is the difference between recovering the clone structure and recovering nothing, at
zero extra cost in phlebotomy.
\item ${\sim}18$ draws per patient, concentrated in a \textbf{burst} bracketing any
qualifying interruption, rather than spread uniformly.
\item \textbf{The window must be measured from drug \emph{clearance}, not from the
last dose.} $u(t)$ is not a step function: an agent with a multi-day half-life is
still exerting selective pressure well after the pill stops. The effective drug-off
window is the nominal hold \emph{minus} the pharmacokinetic washout. Since 21 days is
insufficient and 28 sufficient (\S5.2), the operative requirement is
\textbf{28 days of \emph{cleared} drug} --- for a long-half-life agent this may mean a
nominal hold of 35--40 days, and for a very long-lived agent the design may simply not
be feasible.
\item \textbf{Clone number fitted, not assumed}, with a continuum model in the
candidate set; $\ceff$ and $\aSR$ as the estimands, not per-clone costs.
\item \textbf{No alteration whatsoever to the treatment decision.}
\end{enumerate}

This preserves every property that made the observational framing worth having. No
clinician is asked to accept any evolutionary hypothesis --- one may believe adaptive
therapy is nonsense and still have no objection to the measurement; indeed,
\emph{if} one believes it is nonsense, this is the study that would show it. No
therapeutic risk is introduced. It composes with existing trials as a companion
sampling module. And it is falsifiable.

\begin{center}
\textbf{${\sim}89{,}000$ patients under current practice versus ${\sim}7$--$10$ under
this design, for the same scientific question.}
\end{center}

\section{A falsifiable prediction the design tests on itself}\label{sec:predict}

\S6.4 leaves the design exposed to a confound it cannot fully model. It can, however,
\textbf{detect} it, using data it is already collecting:

\begin{tcolorbox}[keybox]
\textbf{Prediction.} Within the same cohort, $\hat{c}_{\mathrm{eff}}$ estimated from
\emph{toxic} holds will be systematically \emph{lower} than $\hat{c}_{\mathrm{eff}}$
estimated from \emph{non-toxic} holds.

If the strata agree, the host confound is absent and the toxic holds can be pooled.
If they diverge --- in the predicted direction --- the design has detected its own
principal threat from the inside, and the toxic stratum is discarded.
\end{tcolorbox}

This is not a patch. It is the strongest property the design acquired from being
attacked: the confound that could invalidate it is \emph{itself measurable} with the
same eighteen blood draws, because the cohort naturally contains both hold types.

\section{Limitations}\label{sec:limits}

\begin{itemize}[leftmargin=1.4em,itemsep=0.35em]
\item \textbf{Simulation is not evidence about patients.} Everything here is a
statement about information in a model, not a finding about cancer. That is exactly
the claim we intend, and no more.
\item \textbf{This is a cohort design, not a personalised one.} No single patient's
data identifies $\aSR$.
\item \textbf{The ontology is genuinely undecidable, and we route around it rather
than solving it.} A reader who wants to know whether resistance in a given tumour is
clonal will get no help here.
\item \textbf{We did not vary the progression endpoint or the feedback rule.} All
results use a burden threshold and the Gatenby-50\% rule. These are the field's
conventions, not established facts.
\item \textbf{$\ceff$ is coarser than what the models want.} It answers whether the
resistant mass grows more slowly --- enough to decide whether adaptive therapy can
\emph{exist}. Fine-grained \emph{scheduling} would want clonal detail we have shown
to be unidentifiable; \S7.3 suggests this may not matter, since feedback beats
model-based scheduling anyway.
\item \textbf{No immune compartment.} Both model classes are bipartite (sensitive
vs. resistant); immune predation is absent. Note, however, that immune pressure acting
on \emph{total burden} is structurally a host effect of the kind analysed in \S6.4:
benign if it acts symmetrically on the clones, dangerous if it acts
\emph{differentially}. Since immune escape and drug resistance can co-occur, the
differential case is plausible, and a three-compartment audit with an immune predator
is the most important extension we have not done.
\item \textbf{Shedding is assumed to scale linearly with burden.} Necrotic cores and
vascularised margins plausibly shed at different rates, making $\beta$ a function of
tumour size and structure. We tested misspecification of the shedding \emph{ratio}
($\rho$) and the death-coupling ($\kappa$), but not \emph{nonlinearity} in the
burden--shedding map. This is untested and could matter.
\item \textbf{The continuum rescue assumes resistance is smoothly distributed.} If
resistance arises from a single catastrophic macro-evolutionary jump to a distant
phenotype, fitting a continuum could smooth the likelihood surface and yield a
\emph{false positive} for identifiability. We tested discrete truths at $k = 2$ and
$k = 3$; a single, extreme, isolated resistant phenotype is a case we did not test.
\item \textbf{The models are well-mixed; real tumours are spatial.} $\alpha$ is a
mean-field summary of what is physically a local, geometric process. Since
\citet{bacevic2017spatial} find that \emph{spatial} structure constrains resistance,
our well-mixed treatment probably \emph{understates} competition --- but "probably
conservative" is not a substitute for having done the spatial version.
\item \textbf{Parameter sweeps are evidence, not proof.} We varied $c$
(0.05--0.40), assay CV (30--50\%), shedding ratio (0.5--3), $\kappa$ (2--20), clone
number (2--3), competition (0--1.25), and model class (LV, MacArthur). The
conclusions were stable across all of it.
\end{itemize}

\clearpage
\hrule
\vspace{0.6em}
\begin{center}\large\bfseries PART II --- HOW WE GOT HERE\end{center}
\vspace{0.2em}
\hrule

\section{Hypotheses, failures, and revisions}\label{sec:errata}

Part I is what survived. This is what did not.

We report it in full for three reasons: the failures are more instructive than the
successes; a reader is entitled to know how many times the authors were wrong before
arriving at a conclusion they are now asked to trust; and the \emph{pattern} in the
failures is this paper's most transferable finding.

\subsection*{H1 --- ``Oncology samples too slowly. Denser sampling will close the gap.'' --- \textcolor{red!60!black}{REFUTED}}

This was the paper's founding thesis. It was wrong.

The reasoning was seductive: ctDNA clears from plasma with a half-life of minutes to
hours, the dynamics of interest run over days, and restaging happens every eight
weeks. A high-bandwidth sensor read at low bandwidth by convention. Fix the cadence,
fix the problem.

\textbf{What broke it.} We tested it. Six draws $\to$ 281 draws (daily, for forty
weeks). The sign-error rate went from ${\sim}57\%$ to ${\sim}57\%$. \emph{Exactly
unchanged.}

\textbf{What replaced it.} Sampling density is \emph{necessary but nowhere near
sufficient}. The binding constraint was never cadence --- it was that we were
observing an \textbf{unperturbed} system through an observable that \textbf{cannot
see the clones}. Hence the triad (\S5). Density is one leg of three, and worthless on
its own.

\textbf{Why we got it wrong.} We reasoned from an analogy --- signal processing,
Nyquist --- to a system that was not bandwidth-limited but \emph{information}-limited.
The analogy was elegant and irrelevant.

\subsection*{H2 --- ``Clone number is recoverable with the right panel.'' --- \textcolor{red!60!black}{REFUTED}}

Having found that a fixed compartment count destroys the estimate, we proposed
inferring it. We showed a per-variant panel recovers the true clone number 100\% of
the time where a lumped allele fraction recovers it 0\%, and declared the hazard
closed.

\textbf{What broke it.} The candidate set. We had fitted $k = 1, 2, 3, 4$ --- all
discrete. Adding a \textbf{continuum} model, BIC prefers it in \emph{both} worlds,
including 6/8 against a genuinely three-clone truth. We had been asking \emph{which}
discrete model while assuming the answer to \emph{whether}.

\textbf{What replaced it.} The ontology is undecidable, and does not need to be
decided (\S6.3). Re-specify the estimand as a \emph{functional} of the resistance
distribution and the sign comes back correct under either ontology.

\textbf{Why we got it wrong.} We let the candidate set encode the hypothesis. A
model-selection procedure can only tell you which of the models you \emph{offered} it
fits best. Offering only discrete models and concluding ``the tumour is discrete'' is
circular --- and we did it without noticing.

\subsection*{H3 --- ``The cost of resistance is the quantity to measure.'' --- \textcolor{red!60!black}{REFUTED}}

This was not merely our assumption; it is the field's. Adaptive therapy is standardly
justified as working \emph{if and only if} resistance carries a fitness cost.

\textbf{What broke it.} In a world where resistance was entirely \emph{free}
($c = 0$), the model-free adaptive rule still beat MTD substantially. If it works when
resistance costs nothing, the cost is not the mechanism.

\textbf{What replaced it.} $\alpha$ (\S7). Sensitive cells suppress resistant ones by
\textbf{occupying carrying capacity}, not by outgrowing them. The cost amplifies this
but does not create it: it is neither necessary (36 days of benefit at $c \leq 0$) nor
sufficient (4 days at large cost, weak competition), and competition carries
${\sim}2\times$ the leverage. Replicated in a mechanistic consumer--resource model with
a genuine no-competition control.

\textbf{Why we got it wrong.} We inherited the field's framing without auditing it. The
identifiability machinery we had built to \emph{test} the field's parameter is what
eventually revealed the parameter was the wrong one --- but only because we ran a sweep
we had no particular reason to run.

\subsection*{H4 --- ``The benefit is flat in the cost of resistance.'' --- \textcolor{red!60!black}{REFUTED (our own arithmetic)}}

Having discovered H3, we overshot. We reported the benefit as a \textbf{relative} gain,
$(\mathrm{TTP}_{\text{adaptive}} - \mathrm{TTP}_{\text{MTD}})/\mathrm{TTP}_{\text{MTD}}$,
observed it was flat across $c$, and concluded the fitness cost was simply irrelevant.

\textbf{What broke it.} The denominator. MTD survival \emph{also} rises with $c$
(resistant cells grow more slowly, so even MTD does better), masking a real effect in
the numerator. In \textbf{absolute days}, the benefit roughly \textbf{doubles} across
the range of $c$.

\textbf{What replaced it.} The corrected, weaker, true claim: $c$ is not
\emph{required}, but it is not irrelevant either. It roughly doubles the benefit across
its range; $\alpha$ carries twice its leverage; neither is negligible.

\textbf{Why we got it wrong.} We chose a normalisation that flattered a conclusion we
had just become excited about, and did not check whether the denominator was moving.
This is the most ordinary statistical error in this document, and it produced the most
confident-sounding sentence in the draft it appeared in.

\subsection*{H5 --- ``$\aSR$ is structurally unidentifiable.'' --- \textcolor{red!60!black}{REFUTED (our own arithmetic)}}

While auditing H4 we discovered our $\alpha$ results assumed \emph{symmetric}
competition. Relaxing that, the directional coefficient $\aSR$ appeared unrecoverable:
SE $\approx 1.5$, interval spanning zero \emph{and negative} --- the data apparently
could not distinguish suppression from assistance. Extra drug-off windows did not help,
suggesting the barrier was structural. We wrote it up as the paper's most serious
limitation.

\textbf{What broke it.} Two things, both embarrassing.

\begin{enumerate}[leftmargin=1.6em,itemsep=0.3em]
\item Those intervals were \textbf{per-patient}. Every other parameter in this paper is
unresolved per-patient and becomes resolvable by pooling --- $c$ needs $n \approx 7$. We
had silently demanded that $\aSR$ clear a bar in one patient that $c$ had never cleared,
then called the failure structural. Pooled across $n = 10$, \textbf{$\aSR$ resolves
cleanly}: $[0.472, 1.941]$.
\item The SEs came from pseudo-inverting a Fisher matrix with condition number
${\sim}10^{7}$ --- and they were \textbf{non-monotone in the amount of data}
($1.507 \to 1.418 \to 1.553$ as holds went $1 \to 2 \to 3$). \emph{More data cannot
reduce information.} That non-monotonicity was the numerics telling us the bound was
garbage. We read it as a finding.
\end{enumerate}

\textbf{Why we got it wrong.} We committed, twenty minutes later, exactly the error our
own \S4.2 warns against --- trusting a Cram\'er--Rao bound on a near-singular surface.
Writing the warning did not immunise us against the mistake.

\emph{A negative result worth keeping.} We hypothesised that reparameterising to
$\sigma := r_R \aSR$ --- the directly observable slope $-\partial g_R/\partial S$ ---
would be identifiable where $\aSR$ was not, since the data only ever see the product.
\textbf{It buys nothing.} $\sigma$ needs $n \approx 11$, essentially identical. Because
$r_R$ is adequately determined under this design, $\sigma$ and $\aSR$ carry the same
information.

\subsection*{H6 --- ``Vary the resource supply rate to vary competition strength.'' --- \textcolor{orange!70!black}{MIS-SPECIFIED}}

Testing whether $\alpha$-dominance survived outside Lotka--Volterra, our first
mechanistic model varied competition by varying the resource \textbf{supply rate} $\phi$.

\textbf{What broke it.} The $R^*$ rule~\citep{tilman1982resource}. In a
consumer--resource system the equilibrium resource level is
$\mathrm{Res}^* = K m/(r - m)$, which is \textbf{independent of supply}. Raising $\phi$ raises equilibrium \emph{biomass}, not resource
\emph{availability}. Every arm remained fully resource-limited; there was never a
no-competition control.

The mis-specified experiment reported an apparent \textbf{$11\times$ inversion} of the
$\alpha$ result --- a clean-looking, entirely spurious finding.

\textbf{What replaced it.} Niche overlap (\S7.2), which \emph{is} the mechanistic origin
of $\alpha$, and which gives a genuine $\omega = 0$ control. The result then replicated.

\textbf{Why it matters.} This error produced a confident, publishable-looking table
pointing the opposite direction from the truth. Nothing in the output looked wrong. We
caught it only by asking \emph{what does this knob actually do to the equilibrium?} ---
a question we should have asked before running the sweep, not after.


\subsection*{H7 --- ``Any incidental treatment interruption is a free excitation.'' --- \textcolor{red!60!black}{REFUTED (by a referee, correctly)}}

Every draft through v1.0 listed toxicity holds first among the interruptions the design
would exploit. A referee pointed out that this is biologically naive: a Grade 4
neutropenia hold means the patient is in systemic crisis, and the growth and
competition parameters of a tumour in a metabolically deranged host are not those of a
tumour in a well one. We had been treating a physiological catastrophe as a clean
laboratory perturbation.

\textbf{What broke it.} We modelled it (\S6.4). The critique is correct, and the damage
is severe --- but the \emph{structure} of the damage is sharper than the critique
stated. A \textbf{symmetric} host insult is comparatively benign, because $c$ is a
\emph{ratio} and a common factor cancels. A \textbf{differential} insult that spares
the resistant clone is fatal: sign-error goes to \textbf{58\%}, and $\chat$ is biased
\textbf{downward} --- i.e. toward concluding resistance is free, which is exactly the
error the paper exists to prevent. Fitting the host effect as a nuisance parameter only
partly helps (36\% $\to$ 22\%).

\textbf{What replaced it.} An inclusion criterion, not a statistical fix. Only
interruptions that perturb \emph{the drug and nothing else} --- insurance lapse, supply
interruption, scheduled surgery, travel, non-adherence --- are clean excitations.
Toxicity holds are excluded from the identifying set. Enrolment rises by $1/f$; validity
is preserved. And the design gains a way to detect the confound on itself (\S9).

\textbf{Why we got it wrong.} We were reasoning about $u(t)$ and forgot there was a
patient attached to it. The model had a drug term and a tumour term and no host, so the
host could not appear in the analysis --- a category of error no amount of internal
verification would ever have caught, because the machinery was working perfectly on the
wrong system. It took someone outside the model to see it.

\subsection{What survived, and the pattern}

\begin{tcolorbox}[keybox]
\textbf{Every substantive thesis in this paper's first five drafts was refuted. The
design recommendation survived all of them --- and then a referee found the one thing
that could break it.}
\end{tcolorbox}

That asymmetry is worth dwelling on. The \emph{estimand} moved three times (from $c$, to
$\ceff$, to $\aSR$). The \emph{ontology} collapsed. The \emph{mechanism} inverted. But
``clone-resolved observation $+$ a drug-off window $+$ burst sampling around it, at
${\sim}18$ draws and ${\sim}10$ patients, purely observationally'' was right from the
first audit and never moved.

H7 is the exception, and it is the instructive one. It did not come from us. Six
rounds of adversarial self-audit, a verification suite that checks the Fisher
information against the numerical Hessian to $2\times10^{-8}$, and three independent
implementations of the model agreeing to $10^{-6}$ --- none of it could find H7,
because none of it was looking outside the model. The machinery was working perfectly
on a system that was missing a host. \textbf{Internal rigour cannot find a missing
compartment.} Only a reader who knows what a Grade 4 neutropenia hold actually does to
a patient can find that, and no amount of Monte Carlo substitutes for one.

This is what one would hope for from an identifiability analysis. \textbf{Conclusions
about what a protocol can \emph{see} are far more robust than conclusions about what the
world \emph{is}.} That is an argument for doing the audit first.

And the pattern in the failures:

\begin{center}
\small
\begin{tabular}{ll}
\toprule
\textbf{Failure} & \textbf{Category}\\
\midrule
H1 --- sampling density & Reasoned from a wrong analogy (bandwidth, not information)\\
H2 --- clone number & Let the candidate set encode the hypothesis\\
H3 --- cost of resistance & Inherited the field's framing without auditing it\\
H4 --- ``flat in $c$'' & Chose a normalisation that flattered the conclusion\\
H5 --- $\aSR$ & Trusted a bound we had ourselves warned against\\
H6 --- supply rate & Did not ask what the knob controlled\\
H7 --- toxicity holds & Modelled the drug and the tumour, and forgot the host\\
\bottomrule
\end{tabular}
\end{center}

\textbf{Not one of these was a coding error} --- and the most dangerous one was not
even a reasoning error \emph{within} the model. It was a missing piece of biology.

We ran an independent verification suite over the machinery: three model implementations
cross-checked to $<10^{-6}$; the piecewise integrator against brute-force Euler to
$1.3\times10^{-4}$; the Fisher information against the numerical Hessian of the expected
likelihood to $2\times10^{-8}$; the allele-fraction observable checked for saturation;
the headline result checked across independent seed sets. \textbf{All passed.} The code
was right every time.

\begin{tcolorbox}[keybox]
\textbf{Every error was one of interpretation, or of omission.} That is almost certainly
where the errors in this literature live too --- and it is the strongest argument we can
make for why an identifiability audit of \emph{other people's} models is worth doing at
all. It is also an argument for showing your audit to someone who knows the biology, and
believing them when they tell you the perturbation you have been counting on is a sick
patient.
\end{tcolorbox}

\section{Conclusion}

Oncology cannot currently measure the parameters its evolutionary models depend on. Not
the cost of resistance --- roughly half of honest fits get its sign wrong. Not the
competition coefficient --- its estimate wanders the entire admissible range. And not at
any sampling rate, because the system is never perturbed and the observable cannot see
the clones.

The repair is cheap and requires no clinician to accept anything. Resolve the clones.
Sample densely around the drug-free windows that routine care already hands over for
nothing. Ten patients, eighteen draws each, no change to any treatment decision.

And the quantity the field has organised itself around --- the fitness cost of resistance
--- is not the one that decides the question. Competition is. A fitness cost is neither
necessary nor sufficient for adaptive therapy to work; it merely amplifies an effect that
comes from sensitive cells occupying space that resistant cells need. This holds in the
field's own model class, and replicates in a mechanistic one with a genuine control.

We would not want a reader's takeaway to be any single number in this paper;
\S\ref{sec:errata} records how many of them we have had to withdraw. The takeaway is the method:

\begin{tcolorbox}[keybox]
\begin{center}
\textbf{Ask what your measurement protocol can identify,\\before you argue about what your
model implies.}
\end{center}
\end{tcolorbox}

A decade of argument has been conducted over parameters that, under the field's own model
and its own data, could not have been measured either way.

\vspace{1.5em}
\hrule
\vspace{0.6em}

{\small\itshape Code, models, and the full audit --- including every experiment that
refuted an earlier draft --- are reproducible from the accompanying scripts.
Correspondence: ORCID 0009-0006-6283-1696.}

\bibliographystyle{unsrtnat}
\bibliography{references}

\end{document}
